\documentclass{article}
\usepackage[left=2cm, right=2cm, top=2cm, bottom=2cm]{geometry}
\usepackage{graphicx}
\usepackage{amsmath}
\usepackage{amssymb}
\usepackage{amsthm}
\usepackage{fancyhdr}
\usepackage{verbatim}
\usepackage{listings}
\usepackage{xcolor}
\usepackage{pdfpages}
\usepackage{cancel}
\usepackage{physics}

\lstset{
    language=C++,
    basicstyle=\ttfamily\footnotesize,
    keywordstyle=\color{blue},
    commentstyle=\color{green},
    stringstyle=\color{red},
    numbers=left,
    numberstyle=\tiny\color{gray},
    frame=single,
    breaklines=true,
    captionpos=b,
}

\title{MTH 553: Lecture \#22}
\author{Cliff Sun}

\newtheorem{theorem}{Theorem}[section]
\newtheorem{lemma}[theorem]{Lemma}
\newtheorem{definition}[theorem]{Definition}
\newtheorem{conjecture}[theorem]{Conjecture}
\newtheorem{proposition}[theorem]{Proposition}
\newtheorem{corollary}[theorem]{Corollary}

\pagestyle{fancy}
\lhead{\textbf{\thepage}\ \ \nouppercase{\rightmark}}
\chead{MTH 553: Lecture \#22}
\rhead{Cliff Sun}

\begin{document}

\maketitle

\section*{Lecture Span}
\begin{enumerate}
    \item Mean value formula for the heat equation
    \item Strong maximum principle
\end{enumerate}

\section*{Mean Value Formula for the Heat Equation}

\begin{theorem}[Mean Value Theorem]
Let $u \in C^\infty(\Omega_T)$ and assume $u_t = \Delta u$ in $\Omega_T$.  
Let $E(x,t;r) \subset \Omega_T$. Then
\[
u(x,t) = \frac{1}{4r^n}
\int_{E(x,t;r)}
u(y,s)\frac{|x-y|^2}{(t-s)^2}\, dy\, ds .
\]

This is a weighted mean value formula.

\textbf{Interpretation:} The temperature in the heat ball $E(x,t;r)$ determines the temperature at $(x,t)$.
\end{theorem}

\begin{proof}

By translation we may assume $(x,t)=(0,0)$.

\textbf{Step 1:} Assume $u \equiv 1$. Then the right-hand side becomes
\[
\frac{1}{4r^n}\int_{E(0,0;r)} \frac{|y|^2}{s^2}\, dy\, ds .
\]

Introduce the change of variables
\[
y = \xi \sqrt{-s}, \qquad s = -\sigma .
\]

Then
\[
dy = (-s)^{n/2} d\xi, \qquad ds = -d\sigma .
\]

Hence
\[
\int_{E(0,0;r)} \frac{|y|^2}{s^2} dy ds
=
\int \int |\xi|^2 \sigma^{\frac{n}{2}-2} d\xi d\sigma .
\]

Switching to spherical coordinates gives
\[
\int_0^{r^2} \int_0^\infty
\rho^{n+1} e^{-\rho^2/(4\sigma)} d\rho d\sigma .
\]

Using the Gamma function identity
\[
\Gamma(z+1) = z\Gamma(z),
\]
we obtain that the integral equals $4r^n$. Therefore the formula holds for $u\equiv 1$.

\medskip

\textbf{Step 2:} Suppose $u_t = \Delta u$. Define

\[
\varphi(r) =
\frac{1}{4r^n}
\int_{E(0,0;r)}
u(y,s)\frac{|y|^2}{s^2} dy ds .
\]

We show $\varphi'(r)=0$.

Differentiating under the integral sign gives
\[
\varphi'(r)
=
\int_{E(0,0;r)}
\nabla u \cdot \nabla \psi
+
u_t \psi
\, dy ds .
\]

Using $u_t=\Delta u$ we obtain

\[
\varphi'(r)
=
\int_{E(0,0;r)}
\nabla u \cdot \nabla \psi
+
(\Delta u)\psi
\, dy ds .
\]

Let
\[
A = \int_{E} \nabla u \cdot \nabla \psi,
\qquad
B = \int_E (\Delta u)\psi .
\]

Applying integration by parts and the divergence theorem,

\[
B
=
\int_E \nabla \cdot (\psi \nabla u) - \nabla u \cdot \nabla \psi .
\]

Hence

\[
A + B = 0 .
\]

Thus $\varphi'(r)=0$, so $\varphi$ is constant. Taking the limit as $r\to0$ yields

\[
\varphi(r) = u(0,0).
\]

Therefore

\[
u(0,0)
=
\frac{1}{4r^n}
\int_{E(0,0;r)}
u(y,s)\frac{|y|^2}{s^2} dy ds .
\]

\end{proof}

\section*{Strong Maximum Principle for the Heat Equation}

\textbf{Note:} If $u$ satisfies the heat equation, then maxima propagate forward in time.

\begin{theorem}[Strong Maximum Principle]
Let $\Omega \subset \mathbb{R}^n$ be a bounded domain and $u\in C^\infty(\Omega_T)$ satisfy
\[
u_t = \Delta u .
\]

Then
\[
\max_{\Omega_T} u = \max_{\partial_p \Omega_T} u
\]
(where $\partial_p\Omega_T$ is the parabolic boundary).

Furthermore, if for some interior point $(x_0,t_0)\in\Omega_T$ we have

\[
u(x_0,t_0) = \max_{\Omega_T} u,
\]

then $u$ is constant in $\Omega \times (-\infty,t_0]$.
\end{theorem}

\begin{proof}[Sketch]

Assume $t<t_0$ and connect $(x,t)$ to $(x_0,t_0)$ by a path contained in $\Omega_T$.

Cover this path with overlapping heat balls
\[
E_0, E_1, \dots , E_k .
\]

Using the mean value formula,

\[
u(x_0,t_0)
=
\frac{1}{4r^n}
\int_{E_0}
u(y,s)\frac{|x_0-y|^2}{(t_0-s)^2} dy ds .
\]

Since $u(y,s) \le M := u(x_0,t_0)$, equality implies

\[
u(y,s) = M \quad \text{for all } (y,s)\in E_0.
\]

Repeating this argument along the chain of heat balls yields

\[
u(x,t)=M
\quad
\text{for all } x\in\Omega, \; t<t_0 .
\]

Hence $u$ is constant.

\end{proof}

\end{document}